% Cote 119: offprints, reports, courses, CV and letters (typed and printed part): archive cote 119, site page 244 (PDF page 245).
% typescript. Draft compared: none.
% AI-transcribed by Claude Opus 5.5 (claude-opus-5-5, adaptive thinking, effort high), Message Batches API, batch r119_full (msgbatch_01DwZNm5Y7YsKBWyAN4sFs8w), 2026-11-13; not yet checked by eye.
\scanpage{119/244}
\mspage{4}

étudiants comme aux enseignants, par des amphithéâtres surpeuplés, rendant
pratiquement impossible un contact pédagogique vivant entre enseignant et
enseigné.

Du point de vue des étudiants, c'est évidemment c) qui constitue
le vice principal, vice qui a sans doute été une des causes déterminantes
pour créer chez ceux-ci un état d'esprit qui a abouti à la salutaire crise
de l'Université que nous sommes en train de traverser. Alors que pour ce
dernier vice de fonctionnement, c'est dans une large mesure l'insuffisance
des crédits alloués qui semble être en cause, il semble que a) et b) doivent
être imputés surtout à la confusion entre le rôle du maître-enseignant et du
maître-chercheur. En fait, il se trouve que l'importance de la fonction
sociale du premier est généralement méconnue au profit du second, auquel
s'attache traditionnellement un prestige supérieur. Par d'anciennes habitudes
de pensée, le rôle du premier n'est considéré que comme un des attributs du
second. Contre l'enseignant universitaire qui n'est ``\uline{que}'' enseignant, il
y a dans le système actuel une prévention injustifiée et injuste, qui
vis-à-vis des maîtres assistants se traduit par une situation subalterne du
point de vue matériel comme du point de vue participation à la vie de l'Uni-%
versité, et vis-à-vis des maîtres de conférence ou professeurs, à une méses-%
time de la part de nombreux de leurs collègues, qui estiment (avec raison
parfois du point de vue des critères en vigueur jusqu'à maintenant) que
ceux-ci ont réussi à obtenir ``au rabais'' leur titre de docteur et leur
nomination à une chaire.

Le maintien de critères de recrutement et de valeur dépassés, et
l'état d'esprit qu'on vient de mentionner, créent d'ailleurs pour nombre
de scientifiques (et surtout de jeunes scientifiques) des problèmes psycho-%
logiques et des problèmes de conscience tout à fait sérieux. Il y a en effet
le cas de scientifiques de valeur, ayant une culture étendue, un jugement
et un ``goût'' scientifiques très sûrs, des aptitudes pédagogiques excellentes,
qui sont animés d'un enthousiasme sincère pour leur science, mais qui res-%
tent de longues années sans faire oeuvre vraiment originale, et qui n'y
parviennent parfois jamais. Alors que ces hommes seraient précieux pour

\margin{244}
